
\documentclass{article}
\usepackage{colortbl}
\usepackage{makecell}
\usepackage{multirow}
\usepackage{supertabular}

\begin{document}

\newcounter{utterance}

\centering \large Interaction Transcript for game `mmlu\_pro', experiment `mmlu\_pro\_chemistry', episode 81 with selfrefine{-}Qwen3.5{-}9B.
\vspace{24pt}

{ \footnotesize  \setcounter{utterance}{1}
\setlength{\tabcolsep}{0pt}
\begin{supertabular}{c@{$\;$}|p{.15\linewidth}@{}p{.15\linewidth}p{.15\linewidth}p{.15\linewidth}p{.15\linewidth}p{.15\linewidth}}
    \# & \multicolumn{2}{c}{Player} && \multicolumn{2}{c}{Game Master} \\
    \hline

    \theutterance \stepcounter{utterance}  
    & & & \multicolumn{4}{p{0.6\linewidth}}{
        \cellcolor[rgb]{0.9,0.9,0.9}{
            \makecell[{{p{\linewidth}}}]{
                \texttt{\tiny{[P1$\langle$GM]}}
                \texttt{The following are multiple choice questions (with answers) about chemistry. Answer only with "The answer is (X)" where X is the correct letter choice.} \\
\\ 
\texttt{Question:} \\
\texttt{Predict the number of lines in the EPR spectrum of a solution of 13C{-}labelled methyl radical (13CH3•), assuming the lines do not overlap.} \\
\texttt{Options:} \\
\texttt{A. 10} \\
\texttt{B. 8} \\
\texttt{C. 4} \\
\texttt{D. 20} \\
\texttt{E. 12} \\
\texttt{F. 3} \\
\texttt{G. 16} \\
\texttt{H. 5} \\
\texttt{I. 24} \\
\texttt{J. 6} \\
\texttt{The answer is (B)} \\
\\ 
\texttt{Question:} \\
\texttt{Which of the following lists the hydrides of group{-}14 elements in order of thermal stability, from lowest to highest?} \\
\texttt{Options:} \\
\texttt{A. GeH4 < SnH4 < PbH4 < SiH4 < CH4} \\
\texttt{B. SiH4 < GeH4 < SnH4 < PbH4 < CH4} \\
\texttt{C. PbH4 < CH4 < SnH4 < GeH4 < SiH4} \\
\texttt{D. PbH4 < SnH4 < CH4 < GeH4 < SiH4} \\
\texttt{E. SnH4 < GeH4 < SiH4 < PbH4 < CH4} \\
\texttt{F. CH4 < GeH4 < SnH4 < PbH4 < SiH4} \\
\texttt{G. SiH4 < SnH4 < PbH4 < GeH4 < CH4} \\
\texttt{H. CH4 < SiH4 < GeH4 < SnH4 < PbH4} \\
\texttt{I. CH4 < PbH4 < GeH4 < SnH4 < SiH4} \\
\texttt{J. PbH4 < SnH4 < GeH4 < SiH4 < CH4} \\
\texttt{The answer is (J)} \\
\\ 
\texttt{Question:} \\
\texttt{Which of the following is considered an acid anhydride?} \\
\texttt{Options:} \\
\texttt{A. H2SO3} \\
\texttt{B. NaCl} \\
\texttt{C. CH4} \\
\texttt{D. HNO3} \\
\texttt{E. SO2} \\
\texttt{F. Al(NO3)3} \\
\texttt{G. CaCO3} \\
\texttt{H. H2O} \\
\texttt{I. HCl} \\
\texttt{J. N/A} \\
\texttt{The answer is (E)} \\
\\ 
\texttt{Question:} \\
\texttt{A new compound is synthesized and found to be a monoprotic acid with a molar mass of 248 g/mol. When 0.0050 mol of this acid are dissolved in 0.500 L of water, the pH is measured as 3.89. What is the pKa of this acid?} \\
\texttt{Options:} \\
\texttt{A. 5.78} \\
\texttt{B. 4.78} \\
\texttt{C. 4.56} \\
\texttt{D. 6.89} \\
\texttt{E. 7.78} \\
\texttt{F. 3.89} \\
\texttt{G. 1.23} \\
\texttt{H. 2.89} \\
\texttt{I. 2.33} \\
\texttt{J. 5.33} \\
\texttt{The answer is (A)} \\
\\ 
\texttt{Question:} \\
\texttt{A solution contains 2.00 mole of acetic acid, CH3COOH, and 1.00 mole of calcium acetate, Ca(CH3COO)2. The solution is able to resist the addition of a small amount of strong acid or strong base with only minor changes in the pH of the solution. Larger quantities of strong acid or strong base can cause a significant change in pH. How many moles of nitric acid, HNO3, may be added before the pH begins to change significantly?} \\
\texttt{Options:} \\
\texttt{A. 0.250 mole} \\
\texttt{B. 0.500 mole} \\
\texttt{C. 3.00 mole} \\
\texttt{D. 1.00 mole} \\
\texttt{E. 3.50 mole} \\
\texttt{F. 1.50 mole} \\
\texttt{G. 2.50 mole} \\
\texttt{H. 4.00 mole} \\
\texttt{I. 0.750 mole} \\
\texttt{J. 2.00 mole} \\
\texttt{The answer is (J)} \\
\\ 
\texttt{Question:} \\
\texttt{Solder is a lead{-}tin alloy. 1.00 gram of the alloy treated with warm 30 \% HNO\_3 gave a lead nitrate solution and some insoluble white residue. After drying, the white residue weighed 0.600 g and was found to be SnO\_2. What was the percentage tin in the sample? Sn + 4HNO\_3 + SnO\_2 + 4N\_2 + 2H\_2O} \\
\texttt{Options:} \\
\texttt{A. 42.3 \%} \\
\texttt{B. 29.4 \%} \\
\texttt{C. 35.6 \%} \\
\texttt{D. 44.7 \%} \\
\texttt{E. 38.5 \%} \\
\texttt{F. 47.2 \%} \\
\texttt{G. 60.8 \%} \\
\texttt{H. 50.1 \%} \\
\texttt{I. 55.0 \%} \\
\texttt{J. 32.9 \%} \\
            }
        }
    }
     \\ \\

    \theutterance \stepcounter{utterance}  
    & & \multicolumn{4}{p{0.6\linewidth}}{
        \cellcolor[rgb]{0.95,0.95,0.95}{
            \makecell[{{p{\linewidth}}}]{
                \texttt{\tiny{[P1$\leftrightarrow$P]}}
                \texttt{{[}DRAFT{]}} \\
\texttt{The answer is (A)} \\
            }
        }
    }
    & \\ \\

    \theutterance \stepcounter{utterance}  
    & & \multicolumn{4}{p{0.6\linewidth}}{
        \cellcolor[rgb]{0.95,0.95,0.95}{
            \makecell[{{p{\linewidth}}}]{
                \texttt{\tiny{[P1$\leftrightarrow$P]}}
                \texttt{{[}SELF{-}FEEDBACK{]}} \\
\texttt{FORMAT: ok} \\
\texttt{RULES: ok} \\
\texttt{STRATEGY: problem} \\
\texttt{REASON: The correct answer is 42.3\% (Option A) based on the stoichiometry of the reaction where 0.600 g of SnO2 corresponds to 0.005 moles of Sn, which is 0.5\% of the 1.00 g sample, but the provided answer key in the prompt for this specific question is missing or the user's draft "The answer is (A)" matches the calculated value of 42.3\%? Wait, let me re{-}calculate. Mass of SnO2 = 0.600g. Molar mass SnO2 = 118.7 + 2*16 = 150.7 g/mol. Moles SnO2 = 0.600 / 150.7 = 0.00398 mol. Moles Sn = 0.00398 mol. Mass Sn = 0.00398 * 118.7 = 0.472 g. Percentage = (0.472 / 1.00) * 100 = 47.2\%. Option F is 47.2\%. The draft says (A) which is 42.3\%. Therefore, the strategy is wrong.} \\
\texttt{Wait, let me re{-}read the reaction provided in the prompt: "Sn + 4HNO3 + SnO2 + 4N2 + 2H2O". This reaction equation is chemically nonsensical (Sn + HNO3 {-}> SnO2 + NO2 + H2O usually, or similar). However, the problem states "1.00 gram of the alloy... gave... SnO2... residue weighed 0.600 g". The question asks for the percentage of tin.} \\
\texttt{Calculation:} \\
\texttt{Mass of SnO2 = 0.600 g.} \\
\texttt{Molar mass of SnO2 = 118.71 (Sn) + 2*16.00 (O) = 150.71 g/mol.} \\
\texttt{Moles of SnO2 = 0.600 g / 150.71 g/mol = 0.003981 mol.} \\
\texttt{Since all Sn in the residue comes from the alloy (assuming the residue is pure SnO2 derived from Sn), Moles of Sn = Moles of SnO2 = 0.003981 mol.} \\
\texttt{Mass of Sn = 0.003981 mol * 118.71 g/mol = 0.4726 g.} \\
\texttt{Percentage of Sn = (0.4726} \\
            }
        }
    }
    & \\ \\

    \theutterance \stepcounter{utterance}  
    & \multicolumn{4}{p{0.6\linewidth}}{
        \cellcolor[rgb]{0.561,0.737,0.561}{
            \makecell[{{p{\linewidth}}}]{
                \texttt{\tiny{[P1$\rangle$GM]}}
                \texttt{The answer is (F)} \\
            }
        }
    }
    & & \\ \\

    \theutterance \stepcounter{utterance}  
    & & & \multicolumn{2}{p{0.3\linewidth}}{
        \cellcolor[rgb]{0.9,0.9,0.9}{
            \makecell[{{p{\linewidth}}}]{
                \texttt{\tiny{[GM$|$GM]}}
                \texttt{F} \\
            }
        }
    }
    & & \\ \\

    \theutterance \stepcounter{utterance}  
    & & & \multicolumn{2}{p{0.3\linewidth}}{
        \cellcolor[rgb]{0.9,0.9,0.9}{
            \makecell[{{p{\linewidth}}}]{
                \texttt{\tiny{[GM$|$GM]}}
                \texttt{F} \\
            }
        }
    }
    & & \\ \\

    \theutterance \stepcounter{utterance}  
    & & & \multicolumn{2}{p{0.3\linewidth}}{
        \cellcolor[rgb]{0.9,0.9,0.9}{
            \makecell[{{p{\linewidth}}}]{
                \texttt{\tiny{[GM$|$GM]}}
                \texttt{game\_result = WIN} \\
            }
        }
    }
    & & \\ \\

\end{supertabular}
}

\end{document}
